% Generated draft: requires mathematical and source review.
\documentclass[11pt]{article}
\usepackage[margin=25mm]{geometry}
\usepackage{fontspec}
\setmainfont{Latin Modern Roman}
\usepackage{amsmath,amssymb,mathtools,unicode-math}
\setmathfont{Latin Modern Math}
\usepackage{enumitem,needspace,xurl,hyperref}
\setlength{\parindent}{0pt}
\setlength{\parskip}{5pt}
\setlength{\emergencystretch}{3em}
\newcommand{\N}{\mathbb{N}}
\newcommand{\Z}{\mathbb{Z}}
\newcommand{\Q}{\mathbb{Q}}
\newcommand{\R}{\mathbb{R}}
\newcommand{\C}{\mathbb{C}}
\newcommand{\F}{\mathbb{F}}
\newcommand{\A}{\mathbb{A}}
\newcommand{\PP}{\mathbb{P}}
\newcommand{\E}{\mathbb{E}}
\newcommand{\eps}{\varepsilon}
\newcommand{\dd}{\,\mathrm{d}}
\begin{document}
\section*{UnsolvedMath ID: 800254}
\textbf{Source problem code:} catalogue entry
\subsection{Definitions and mathematical statement}
Under the usual several-complex-variables interpretation, let \(n\geq 1\) be an integer and let \(\Omega\subseteq\C^n\) be a domain, meaning a nonempty connected open set. Its holomorphic automorphism group is
\[
\operatorname{Aut}(\Omega)=\{f:\Omega\to\Omega: f\text{ is bijective and both }f\text{ and }f^{-1}\text{ are holomorphic}\},
\]
with composition as the group operation. Equip this group with the compact-open topology, that is, the topology of uniform convergence on compact subsets of \(\Omega\).

A standard interpretation of smoothness is that \(\Omega\) has \(C^\infty\) boundary: for every \(p\in\partial\Omega\), there exist a neighborhood \(U\) of \(p\) and a real-valued function \(\rho\in C^\infty(U)\) such that
\[
\Omega\cap U=\{z\in U:\rho(z)<0\},\qquad
\partial\Omega\cap U=\{z\in U:\rho(z)=0\},\qquad
\dd\rho\neq 0\text{ on }\partial\Omega\cap U.
\]

\textbf{Classification question (source Problem 53).} Which smooth domains \(\Omega\) have noncompact automorphism group \(\operatorname{Aut}(\Omega)\)?

The definitions above make the customary terminology explicit, but the supplied statement does not specify the ambient dimension, the intended degree of boundary regularity, or the topology on the automorphism group. It also does not say whether domains must be bounded or satisfy additional geometric hypotheses. Thus these conventions are an interpretation, and the intended class of domains remains unresolved; boundedness and further geometric assumptions are not supplied hypotheses.
\subsection{Short English statement}
Characterize smooth domains whose group of biholomorphic self-maps is noncompact. The supplied source leaves the dimension and intended class of domains unspecified; the boundary regularity and group topology used above are conventional interpretations requiring review.
\subsection{Sources}
\begin{itemize}
\item Local UnsolvedMath registry, supplied complete record: permanent ID \(800254\), exact source problem code \texttt{catalogue entry}, title \emph{Several Complex Variables Problem List --- Problem 53}. The registry attributes the problem to Buzzard, Fornaess, and Krantz, with proposed year 1995. Its status label is \texttt{open}; this is supplied source metadata, not an independently verified mathematical assessment.
\item Buzzard, Fornaess, and Krantz, \emph{The Several Complex Variables Problem List} (1995), Problem 53. Supplied source URL: \url{https://arxiv.org/abs/math/9509202}. The record reports access on 2026-10-16 and extraction method \texttt{arxiv\_latex}. These provenance details are supplied metadata; the source has not been independently visited or verified for this draft.
\end{itemize}
\end{document}
